\section{Protocol-Aware Comparison With Prior Work}

Table~\ref{tab:literature} is an estimand audit, not a classifier leaderboard.
An internal result trains and tests within one resource; a chronological result
changes calendar period; and a frozen external result applies a source-selected
model to a different resource. We use three comparison classes: A-close for a
substantially aligned endpoint, modality, and evaluation direction; B-partial
when at least one major protocol axis differs; and C-context when direct metric
ranking would be misleading. Even A-close rows are not formal head-to-head
tests because cohort releases and implementation details differ.

\begingroup
\scriptsize
\setlength{\tabcolsep}{3pt}
\begin{longtable}{L{0.15\textwidth}L{0.14\textwidth}L{0.17\textwidth}L{0.12\textwidth}L{0.15\textwidth}L{0.17\textwidth}}
\caption{Full-text-audited respiratory-audio literature, organized by the
quantity each experiment estimates. Target use includes fitting, cohort
selection, representation selection, thresholding, or calibration.}
\label{tab:literature}\\
\toprule
Study & Data and inputs & Unit and protocol & Reported result & Target-data use & Comparison boundary \\
\midrule
\endfirsthead
\toprule
Study & Data and inputs & Unit and protocol & Reported result & Target-data use & Comparison boundary \\
\midrule
\endhead
Bhattacharya et al., \emph{Coswara} \cite{bhattacharya2023coswara} & Coswara; nine audio tasks plus symptoms & Participant; 70/15/15 internal split & AUROC 0.915 (95\% CI 0.885--0.941) & Internal validation and test & B-partial: different snapshot, tasks, symptoms, and model \\
Orlandic et al., \emph{COUGHVID} \cite{orlandic2021coughvid} & COUGHVID cough events & Recording; nested 10-fold internal CV & AUC $0.964\pm0.033$ & Internal & C-context: cough-versus-noncough detection, not COVID classification \\
Xia et al., \emph{COVID-19 Sounds} \cite{xia2021covid} & Cough, breath, and voice & Participant; speaker-independent 7:1:2 split & Multimodal AUROC 0.71 (95\% CI 0.65--0.76) & Internal & B-partial: multimodal but a different cohort and label process \\
Chetupalli et al. \cite{chetupalli2023multimodal} & Coswara deep breath, heavy cough, counting, symptoms & Participant; subject-disjoint 80/20 with development CV & Audio-only 0.88; audio+symptoms 0.963 & Internal & A-close internal context for 0.88; symptom-assisted 0.963 is C-context \\
Truong et al., FAIR \cite{truong2024fair} & Coswara; seven sound instances & Participant; fixed test plus rotating development folds and two seeds & AUROC $0.8658\pm0.0115$ & Internal & A-close internal context; different cohort and seven-input architecture \\
Han et al. \cite{han2022realistic} & COVID-19 Sounds; cough, breath, voice & Participant; realistic and deliberately biased internal designs & Realistic multimodal 0.71 (95\% CI 0.65--0.77); biased designs up to 0.90 & Internal & B-partial: direct evidence that cohort construction changes AUROC \\
Coppock et al. \cite{coppock2024audio} & 67,842 PCR-linked UK participants; cough, exhalation, sentence & Participant; random, matched, longitudinal, and longitudinal-matched & Sentence 0.846 random versus 0.619 matched & Internal, prespecified covariate controls & B-partial: strongest methodological comparator; substantially larger PCR cohort \\
Ganitidis et al. \cite{ganitidis2025drift} & Coswara and COVID-19 Sounds cough & Participant-separated chronological 70/30 plus adaptation experiments & Coswara 0.668 to 0.597; COVID-19 Sounds 0.691 to 0.607 & Later labels used for evaluation and adaptation & B-partial: close temporal question, but cough-only and adaptation-focused \\
Kim and Lee \cite{kim2024longitudinal} & Cambridge source; Virufy and Coswara periods & Cough; dataset and period transfer; participant grouping not explicit & Cambridge 0.9346; Coswara 0.8250, 0.7724, 0.5509 by period & No target retraining reported & B-partial: acquisition, calendar, presumed variant, and prevalence change together \\
Grant et al. \cite{grant2022deployment} & DiCOVA/Coswara source; device and source targets & Recording/participant as available; source CV and stress tests & Web cough 0.42 versus Android cough 0.75; mean 0.77 & Frozen and target-optimized thresholds both reported & B-partial: deployment-mismatch evidence; different source-target pairs \\
Atmaja et al. \cite{atmaja2025crossdataset} & Pooled Coswara and COUGHVID training; ComParE test & Cough segments; fixed development and 154-segment test & Unweighted accuracy 0.8819 & The same test supported choices of split, segmentation, augmentation, and mixup & C-context numerically: pooled-source training and test-guided configuration; result is not AUROC \\
Aytekin et al., HST \cite{aytekin2024hst} & Cambridge tasks and separate balanced COUGHVID task & Dataset-specific repeated participant-disjoint internal partitions & COUGHVID AUROC $0.90\pm0.01$ & COUGHVID used for training and selection & C-context for transfer; valid internal architecture comparison only \\
Islam et al., DNDT/DNDF \cite{islam2026robust} & Five cough datasets; selected balanced cohorts & Stratified internal CV and explicit train-one/test-another transfer & Internal: Coswara 0.92, COUGHVID 0.93; cross-data: 0.53 & Target cohort balanced by label; no target fitting in transfer & A-close for Coswara-to-COUGHVID direction and cough modality; target is 680/680 and participant grouping is unreported \\
Pahar et al. \cite{pahar2022transfer} & Coswara; Sarcos target; separate cough, breath, speech & Subject-aggregated nested leave-$p$-out; external target has 44 subjects & Coswara cough 0.982; Coswara$\rightarrow$Sarcos 0.954 & No target fitting reported & B-partial: genuine transfer, but different and very small target \\
Avila et al. \cite{avila2021feature} & DiCOVA/Coswara cough & Fixed participant folds and 233-recording blind test & Blind fusion 0.808; nonnested feature CV 0.96 fell to 0.768 when nested & Blind test excluded from selection & B-partial and direct precedent for feature-selection optimism \\
Laguarta et al. \cite{laguarta2020cough} & Private 5,320-subject forced-cough cohort & Subject; random 80/20 internal split & AUROC 0.97 on the official-test subset & Internal & C-context: private internal task with heterogeneous label ascertainment \\
Chowdhury et al. \cite{chowdhury2022mcdm} & Several cough datasets, separate or merged & Stratified 10-fold CV; participant grouping unreported & Coswara about 0.64--0.66; other internal rows up to 0.97--0.98 & Internal or pooled & C-context for transfer; no frozen train-one/test-another estimate \\
Celik, \emph{CovidCoughNet} \cite{celik2023covidcoughnet} & COUGHVID and Coswara; modalities modeled separately & Recording-level 80/20 split; participant and augmentation-parent grouping unreported & COUGHVID binary AUC 0.9505; Coswara cough 0.9848 & Test set used across pitch-shift comparisons & C-context: merged symptomatic/COVID endpoint, no multimodal fusion or external transfer \\
Hussain et al., \emph{Cough2COVID-19} \cite{hussain2024cough2covid} & Coswara, Virufy, ComParE development; selected COUGHVID target & Reported multi-source leave-one-dataset-out cough evaluation & COUGHVID AUROC 0.981 & Labeled results from all four resources informed representation-family ranking & B-partial: important counterexample; multi-source training, selected 651/660 target, and no uncertainty \\
de Brito et al. \cite{debrito2025audiomae} & Coswara and COUGHVID cough & Stratified 80/20 and bidirectional transfer after demographic balancing & Cross-data AUROC 0.51--0.68 across Audio-MAE/PANN & Target labels define balanced evaluation cohorts & A-close scientifically for Coswara-to-COUGHVID transfer, but non-peer-reviewed and not the full target distribution \\
Lin et al., \emph{SympCoughNet} \cite{lin2025sympcoughnet} & UK cough plus 14 symptoms & Sample; stratified 70/15/15; participant separation not reported & Symptom+audio 0.9474; audio-only ablation 0.8127 & Internal & C-context for headline; B-partial for audio-only ablation \\
\midrule
Present study & Coswara cough, breath, speech; full known-label COUGHVID cough target & Source participant; target recording; internal, chronological, and frozen transfer & Internal multimodal 0.897; matched cough external 0.523--0.543 & No target labels in source model or threshold; separate recalibration sensitivity disclosed & Adds common audit controls, but does not externally validate multimodal fusion \\
\bottomrule
\end{longtable}
\endgroup

The closest internal multimodal context is Chetupalli's audio-only fusion at
0.88 and FAIR at $0.8658\pm0.0115$; the present selected internal estimate is
0.897. These values are contextual, not superiority tests. The closest
peer-reviewed frozen direction is DNDT/DNDF's Coswara-to-COUGHVID AUROC of
0.53, consistent with the present conventional range of 0.523--0.543. The
Audio-MAE/PANN technical report independently places deep transfer in the
0.51--0.68 range, while the present WavLM and CNN--BiGRU estimates are 0.484
and 0.548.

Cough2COVID-19 is a material counterexample rather than a result to dismiss. It
reports 0.981 for multi-source development followed by a selected balanced
COUGHVID target. That estimand differs from single-source Coswara transfer to
all 8,331 known-label COUGHVID recordings, and labeled results from all four
resources informed its representation-family ranking. The contrast therefore
motivates a future controlled factorial study of source diversity, target
selection, and representation choice; it does not establish that either
reported estimate is erroneous.
